% ===== §7 Protecting the Unsymbolisable: The Political Economy of Conditions =====
\section{Protecting the Unsymbolisable: The Political Economy of Conditions}
\label{sec:conditions}

\noindent
The rupture identified at the close of the last section---that the dyad's security is a means to an end the means cannot touch---is not a difficulty to be managed but the central methodological fact of the entire paper, and this section states it in full. It is the section on which the paper's distinctiveness rests. Every theory the paper recruits to the second question---how love is not destroyed by power---is a theory of the means; and the burden here is to establish, exactly, what the means can and cannot do for the end. The answer is severe, and it governs all the practical chapters that follow: \emph{the tools of power and of wealth can protect the conditions under which love occurs, but they cannot produce, secure, or even touch the love itself.}

\subsection{The end is a process of the real}
\label{sub:conditions-real}

The reason lies in what the end \emph{is}. The dynamics of love---the generativity of the coupling, the thing all the defending is for---is a process of the \emph{real}, in the register the series has used throughout. It circles a vacancy, \emph{das Ding}, that cannot be fully taken up into the symbolic; it is not a quantity and admits no measure; it is not a good and enters no exchange. The series established this elsewhere as a thesis about value: that the value generated in a good coupling is non-appropriable and cannot be siphoned, precisely because it is not a thing held but a phase accumulated around a vacancy that no party occupies.\footnote{On the real value that circles the vacancy and cannot be appropriated or siphoned---as against the imaginary value organised around a possessable object---see \citet{paperix}.} What is true of the value is true of the dynamics that generate it: they belong to the real, and the real is, by definition, what escapes complete symbolisation.

\subsection{The instruments are symbolic, and the touch decoheres}
\label{sub:conditions-symbolic}

The instruments of protection, by contrast, are without exception \emph{symbolic}. An IR strategy, a balance, an alliance, a boundary-work, a deployment of relational wealth---each operates by representation, calculation, comparison, the assignment of value to commensurable objects. They live in the symbolic order, which is their power and their limit. For a symbolic instrument applied \emph{directly} to a process of the real does not protect that process; it transforms it, and the transformation is destruction.

\begin{claim}[The decohering touch]
A symbolic instrument cannot act directly upon a process of the real without changing its nature, and the change is loss. Love managed---made the object of strategy, calculation, or deployment---decoheres: counted, it is no longer the uncountable thing it was; secured by guarantee, it is no longer the unsecured thing whose very lack of guarantee was its substance; deployed as wealth, it becomes wealth, and ceases to be love. The attempt to seize the real with a symbolic tool does not capture the real; it replaces it with its symbolic shell. This is the mechanism of the decoherence of trust catalogued in \xref{sub:symptom-three}, and its limiting case is the forged trust of \citet{paperxiii}: the apparatus perfect, the substance gone.
\end{claim}

\noindent
The point is not that managing love is difficult or distasteful. It is that it is a \emph{category error} whose result is the annihilation of the thing managed. One does not protect a flower by handling the bloom; the handling bruises it. This is not a counsel of passivity---there is a great deal to be done---but a specification of \emph{where} the doing must be directed.

\subsection{Fence and bloom}
\label{sub:conditions-fence}

The governing image of the paper's method can now be stated. What the symbolic instruments protect is not the bloom but the \emph{conditions} of the bloom---the soil, the light, the absence of what would trample it. One builds a fence not to cause the flowering, which the fence cannot do, but to keep out the feet and the frost that would prevent it. The fence is wholly in the order of means---it can be designed, built, measured, repaired---and it is wholly necessary, for without it the bloom is at the mercy of the field. But the fence does not bloom, and no perfection of fence-building ever produced a flower. The blooming follows a logic the fence cannot reach: the logic of the real, of \emph{das Ding}, of the desire that circles it.

\begin{claim}[Fence and bloom]
The operations of power can protect only the \emph{surrounding conditions} under which the dynamics of love may occur---the cleared space, the warded boundary, the channel kept clear enough for error to remain attributable---never the dynamics themselves. The fence is the whole domain of strategy: necessary, buildable, and powerless to cause the flowering it makes room for. To confuse the two---to believe that a sufficiently well-built fence is a flower, that perfected conditions \emph{are} love---is the characteristic error of every attempt to engineer intimacy, and it ends, reliably, in a walled and barren plot.
\end{claim}

\subsection{The political economy of conditions, not of love}
\label{sub:conditions-pe}

From this follows the exact and limited title under which IR and the political economy of relational wealth may enter the argument. They are admitted as a \emph{political economy of conditions}---never as a political economy of love. Everything the series has said about relational wealth, everything this paper will say about balancing, hedging, alliance, and boundary-work, belongs to the management of the fence: it concerns the security, the clearing, the warding, the keeping-open of the space in which a coupling may generate. None of it concerns, because none of it \emph{can} concern, the generation itself. The generativity of love is not a managed output of the political economy; it is what the political economy exists to make room for, and falls silent before.

This is a real limit on the paper's own apparatus, and it is stated as such. The holonomic value-theory, the predictive-coding account, the IR strategies---powerful as they are at the level of conditions---are mute at the level of the end. They can tell us when value has failed to accumulate and diagnose the obstacle; they cannot tell us how to make the loop close, for the closing is an event of the real. The most a political economy of conditions can do is clear the ground and keep it clear, and then wait, as a gardener waits, on a process it serves but does not command.

\subsection{Trust as the interface}
\label{sub:conditions-trust}

One element of the system, and one only, crosses the line between the two levels---and this is why \emph{trust} stands in the paper's title beside power and estrangement. Trust is neither wholly a condition nor wholly the bloom. It is the interface between them.

On one side, trust can be \emph{cultivated by symbolic means}, and the series has shown how: by costly and credible signals, by the structural priors a shared history lays down, by the slow projection of a real fidelity into observable, readable form. To this extent trust belongs to the fence---it can be built, tended, and on occasion repaired by instruments of the symbolic order. On the other side, trust is \emph{rooted in the real}: at its limit it is not an inference from evidence but a coupling that cannot be faked, a fidelity that is what it is beneath all its signs, and that the signs at best express and at worst forge.\footnote{On relational trust as spanning the structural, self-exposure, and active-inference channels---and thereby crossing from the symbolic into the real---see \citet{paperxiii}.} Trust is thus the one place where the work done on the conditions reaches \emph{toward} the bloom without ever seizing it: a credible signal does not produce the real coupling, but it clears the ground on which the real coupling may, or may not, occur. Trust is the bridge by which the political economy of conditions touches the hem of what it protects.

\begin{claim}[Trust as interface]
Trust is the sole element of the bond that crosses both levels: cultivable, in part, by symbolic means, and yet rooted in a real coupling that cannot be faked. It is therefore the interface through which the political economy of conditions reaches toward---without ever seizing---the dynamics of love it exists to protect. And it is, as the next chapters will show, the variable that decides whether an asymmetry of power becomes generative or frozen: for an asymmetry suffused with trust empowers, while the same asymmetry voided of trust dominates. Trust runs through every question the paper asks because it is the one thread that runs through both of the levels the paper holds apart.
\end{claim}

\subsection{What this licenses, and what it forbids}
\label{sub:conditions-license}

The section closes by fixing what its central distinction permits and prohibits in the chapters to come. It \emph{licenses} a full and unembarrassed political economy of the fence: the next section's inventory of IR strategies, the praxis of boundary-work and counter-force, the deployment of relational wealth to secure the conditions of a coupling---all of this is legitimate, necessary, and the proper business of the means. It \emph{forbids} the migration of any of these instruments across the interface to the management of the end: the moment a strategy, a balance, or a calculation is turned upon the love itself rather than upon its conditions, it ceases to protect and begins to decohere. The discipline the paper imposes on itself, and recommends, is therefore not a quietism but a \emph{vigilance about direction}---a constant attention to whether a given operation is tending the fence or handling the bloom. With that discipline fixed, the paper can proceed to its inventory of the fence-tending strategies, and to the filter that decides which of them a bond may use without, in the using, trampling what it meant to protect.
